\section{Task Behavior Graphs}
\label{sec:task-behavior-graphs}

The cost model defines typed dimensions and estimate objects. Influence domains
define semantic roles and permitted role flow. A Task Behavior Graph connects
these roles for a task kind or task family.

For task kind $k$, a Task Behavior Graph is a typed directed graph
\[
\mathcal{G}_k
=
\left(
V_k,
E_k,
\delta_k,
\tau,
\lambda,
w,
\gamma
\right).
\]
Here $V_k$ is the set of nodes, $E_k\subseteq V_k\times V_k$ is the set of
directed influence edges, $\delta_k:V_k\rightarrow\mathcal{D}$ assigns an
influence domain to each node, $\tau$ assigns a node type, $\lambda$ assigns an
edge role, $w$ stores edge parameters, and $\gamma$ stores confidence values.

Nodes represent typed model variables rather than raw metric names. A node may
represent task identity, an intrinsic feature, an execution-context feature, a
resource-demand estimate object, a worker-occupancy estimate object, a symptom,
an outcome, a reliability diagnostic, an operational-risk output, or a
profiling-policy trigger. Each node has one primary influence domain; if one raw
observation supports several roles, the mapping layer creates several variables
before they appear in the graph.

Edges represent modeled influence assumptions and must respect the permitted
domain-flow relation from Section~\ref{sec:influence-domain-flow}. Edge roles
include baseline influence, contextual amplification, overhead contribution,
evidence link, symptom link, outcome link, diagnostic link, risk propagation,
and policy trigger:
\[
\lambda(e)
\in
\{
\mathrm{baseline},
\mathrm{amplification},
\mathrm{overhead},
\mathrm{evidence},
\mathrm{symptom},
\mathrm{outcome},
\mathrm{diagnostic},
\mathrm{risk},
\mathrm{policy}
\}.
\]
An edge is not proof of causality. It records an explicit assumption used by the
task-kind model. The confidence value $\gamma_e$ records how reliable that
assumption is considered by the current model state or manifest.

A Task Behavior Graph may implement a dimension-specific estimator $F_{k,q}$,
parameterize it, or document the influence structure used by a separate
estimator. The required property is that the graph makes task-kind-specific
influence assumptions explicit while preserving the role boundaries defined by
influence domains. In particular, outcome, diagnostic, and risk nodes must not be
used as pre-execution intrinsic or context inputs for the same attempt.

A replication graph fragment illustrates the structure. The feature
\texttt{payload\_size\_bytes} may be an intrinsic node that influences
\texttt{network\_output\_bytes} through a baseline edge. The context feature
\texttt{peer\_health} may influence \texttt{worker\_slot\_time} through an
amplification edge. The context feature \texttt{retry\_attempt} may influence
\texttt{connection\_hold\_time} and a retry diagnostic. The cost node
\texttt{worker\_slot\_time} may feed a queue-starvation risk node, and that risk
node may trigger a tail-aware profiling policy.
